\section{Durchführung}
Im ersten Versuchsteil wird die Abhängigkeit des kritischen Stroms $I_\mathrm{c}$ vom magnetischen Feld $B$ untersucht. Hierzu wird an einer Spule Strom angelegt um ein Magnetfeld zu erzeugen. Bei verschiedenen angelegten Strömen werden dann Strom-Spannungs-Kennlinien von \SI{-0.2}{\milli\ampere} bis \SI{0.2}{\milli\ampere} aufgenommen. In der Mitte der Kennlinie ist ein Sprung der Höhe $2 I_\mathrm{c}$ zu sehen.

Im zweiten Versuchsteil wird die Temperaturabhängigkeit des kritischen Stroms $I_\mathrm{c}$ und der Energielücke $\Delta$ untersucht. Hierzu werden an einem Heizelement verschiedenen Ströme angelegt, die Temperatur gemessen und die Strom-Spannungs-Kennlinien von \SI{-1}{\milli\ampere} bis \SI{1}{\milli\ampere} aufgenommen. An diesen Kennlinien sind insgesamt drei Sprünge zu sehen. Einer befindet sich am Ursprung und hat eine Höhe von $2 I_\mathrm{c}$. Die anderen beiden Sprünge befinden sich bei $V = \pm 2\Delta/e$.